\documentclass{article}
\usepackage[english]{babel}
\usepackage{geometry,amsmath,amssymb,enumerate,hyperref,latexsym,theorem}
\geometry{letterpaper}

%%%%%%%%%% Start TeXmacs macros
\newcommand{\assign}{:=}
\newcommand{\tmem}[1]{{\em #1\/}}
\newcommand{\tmmathbf}[1]{\ensuremath{\boldsymbol{#1}}}
\newcommand{\tmscript}[1]{\text{\scriptsize{$#1$}}}
\newenvironment{enumeratenumeric}{\begin{enumerate}[1.] }{\end{enumerate}}
\newenvironment{proof}{\noindent\textbf{Proof\ }}{\hspace*{\fill}$\Box$\medskip}
\newtheorem{definition}{Definition}
\newtheorem{lemma}{Lemma}
\newtheorem{proposition}{Proposition}
{\theorembodyfont{\rmfamily}\newtheorem{remark}{Remark}}
\newtheorem{theorem}{Theorem}
%%%%%%%%%% End TeXmacs macros

\newtheorem{assumption}{Assumption}
\newcommand{\re}{Re}
\newcommand{\R}{\ensuremath{\mathbb{R}}}
\newcommand{\C}{\mathbb{C}}
\newcommand{\N}{\mathbb{N}}
\newcommand{\dd}{\hspace{0.17em} \mathrm{d}}
\newcommand{\ii}{\mathrm{i}}

\begin{document}

\title{The Pad{\'e}--M{\"u}ntz Method forConstant-Coefficient Fractional
Riccati Equations}

\author{Stephen Crowley}

\date{April 2026}

\maketitle

\begin{abstract}
  A purely algebraic, discretization-free solution method is developed for the
  general constant-coefficient fractional Riccati equation $D^{\mu} y = c_0 +
  c_1  \hspace{0.17em} y + c_2  \hspace{0.17em} y^2$ of Caputo order $\mu \in
  (0, 1)$ with initial condition $I^{1 - \mu} y (0) = 0$.\footnote{A reference
  implementation is available in
  \href{https://github.com/crowlogic/arb4j/blob/master/src/main/java/arb/equations/ConstantCoefficientFractionalRiccatiEquation.java}{arb4j}.}
  Because the M{\"u}ntz lattice $\{t^{k \mu} \}_{k \ge 1}$ is closed under the
  Caputo derivative, fractional integration, and pointwise multiplication, the
  M{\"u}ntz spectral Tau substitution yields an explicit Puiseux expansion $y
  (t) = \sum_{k \ge 1} a_k  \hspace{0.17em} t^{k \mu}$ whose coefficients
  satisfy a closed Gamma-ratio convolution recurrence. The induced power
  series in the variable $z = t^{\mu}$ has a strictly positive radius of
  convergence $R$, established by a Banach fixed-point argument on the
  equivalent Volterra equation, and is finite for generic parameters because
  the complexified Riccati flow develops movable singularities in the complex
  $z$-plane; the recurrence by itself therefore produces only a local
  representation and is not globally convergent in $t$. Global-in-$t$
  convergence on the positive real axis is recovered by re-summing the Puiseux
  series via the diagonal $[M / M]$ Pad{\'e} approximant in $z$, constructed
  from the M{\"u}ntz coefficients through a single Hankel linear system $H_M
  \tmmathbf{q}= -\tmmathbf{b}$. Under the de~Montessus de~Ballore and
  Nuttall--Pommerenke theorems, the diagonal Pad{\'e} sequence converges to
  the analytic continuation of $y$ along $z = t^{\mu}$ for all $t \ge 0$
  whenever the solution has no real-positive poles {\tmem{and}} $g (z) \assign
  y (z^{1 / \mu})$ admits a meromorphic continuation to $\C$; in that setting
  a computable a-posteriori error bound $|y (t) - y_M (t) | \le | \Delta_M
  (t^{\mu}) |^2 / (| \Delta_{M - 1} (t^{\mu}) | - | \Delta_M (t^{\mu}) |)$ in
  terms of successive Pad{\'e} differences $\Delta_k = R_k - R_{k - 1}$ is
  derived from a uniform geometric tail estimate. A concluding section
  identifies the deeper structure: the substitution $z = t^{\mu}$ promotes the
  solution to a {\tmem{resurgent transseries}} in $z$, the Pad{\'e} step is
  the computational instantiation of {\tmem{Borel--Pad{\'e} resummation}}, and
  the global analytic continuation is governed by a scalar
  {\tmem{Riemann--Hilbert problem}} whose jump contours are the Stokes lines
  of the transseries. This perspective points to natural generalizations to
  higher-degree polynomial nonlinearities, fractional ODE systems via
  Hermite--Pad{\'e} simultaneous approximation, and non-constant coefficients
  via isomonodromic Riemann--Hilbert problems of Painlev{\'e} type.
\end{abstract}

{\tableofcontents}

\section{Fractional Calculus and the M{\"u}ntz Basis}\label{sec:prelim}

\subsection{Riemann--Liouville integral and Caputo derivative}

\begin{definition}
  [Riemann--Liouville integral]\label{def:RL} Let $r > 0$ and $f \in
  L^1_{\mathrm{loc}} [0, T]$. Define
  \begin{equation}
    \label{eq:RL_int} I^r f (t) \assign \frac{1}{\Gamma (r)}  \int_0^t (t -
    s)^{r - 1} f (s) \hspace{0.17em} \dd s, \qquad t \in [0, T]
  \end{equation}
\end{definition}

\begin{definition}
  [Caputo derivative]\label{def:Caputo} Let $\mu \in (0, 1)$ and $f$ satisfy
  $I^{1 - \mu} f \in C^1 [0, T]$. Define
  \begin{equation}
    \label{eq:caputo} D^{\mu} f (t) \assign \frac{\dd}{\dd t} 
    \hspace{-0.17em} (I^{1 - \mu} f (t)) - \frac{f (0)}{\Gamma (1 - \mu)} t^{-
    \mu}, \qquad t \in (0, T]
  \end{equation}
  The fundamental relations are
  \begin{equation}
    \label{eq:caputo_RL} I^{\mu} D^{\mu} f (t) = f (t) - f (0), \qquad D^{\mu}
    I^{\mu} f (t) = f (t)
  \end{equation}
\end{definition}

\begin{remark}
  [Vanishing correction term for the Riccati solution]\label{rem:correction}
  For the fractional Riccati equation \eqref{eq:frac_riccati} below, the
  initial condition $I^{1 - \mu} y (0) = 0$ is consistent with the M{\"u}ntz
  ansatz $y (t) = \sum_{k \ge 1} a_k  \hspace{0.17em} t^{k \mu}$, since by
  \eqref{eq:Ir_power} below $I^{1 - \mu} [t^{k \mu}] |_{t = 0} = 0$ for all $k
  \ge 1$. The same ansatz gives $y (0) = 0$, so the correction term $- f (0)
  t^{- \mu} / \Gamma (1 - \mu)$ in \eqref{eq:caputo} vanishes identically
  along the Riccati solution, and the Caputo derivative reduces to $D^{\mu} y
  (t) = \frac{\dd}{\dd t} I^{1 - \mu} y (t)$ throughout.
\end{remark}

\begin{proposition}
  [Action on powers]\label{prop:powers} Let $s > - 1$ and $r > 0$. Then
  \begin{equation}
    \label{eq:Ir_power} I^r t^s = \frac{\Gamma (s + 1)}{\Gamma (s + r + 1)} 
    \hspace{0.17em} t^{s + r}
  \end{equation}
  Let $\mu \in (0, 1)$ and $s > 0$. Then
  \begin{equation}
    \label{eq:Dmu_power} D^{\mu} t^s = \frac{\Gamma (s + 1)}{\Gamma (s + 1 -
    \mu)}  \hspace{0.17em} t^{s - \mu}
  \end{equation}
\end{proposition}

\begin{proof}
  For \eqref{eq:Ir_power}, $I^r t^s = \frac{1}{\Gamma (r)}  \int_0^t (t -
  u)^{r - 1} u^s \hspace{0.17em} \dd u$. The substitution $u = tx$ gives
  \[ I^r t^s = \frac{t^{s + r}}{\Gamma (r)}  \int_0^1 (1 - x)^{r - 1} x^s
     \hspace{0.17em} \dd x = \frac{t^{s + r}}{\Gamma (r)} \cdot \frac{\Gamma
     (s + 1)  \hspace{0.17em} \Gamma (r)}{\Gamma (s + r + 1)} = \frac{\Gamma
     (s + 1)}{\Gamma (s + r + 1)}  \hspace{0.17em} t^{s + r} . \]
  For \eqref{eq:Dmu_power}: since $s > 0$, $\lim_{t \to 0^+} t^s = 0$; the
  Caputo correction term $- f (0)  \hspace{0.17em} t^{- \mu} / \Gamma (1 -
  \mu)$ in \eqref{eq:caputo} therefore vanishes, and $D^{\mu} t^s =
  \frac{\dd}{\dd t} I^{1 - \mu} t^s$. By \eqref{eq:Ir_power} with $r = 1 -
  \mu$, $I^{1 - \mu} t^s = \frac{\Gamma (s + 1)}{\Gamma (s + 2 - \mu)} t^{s +
  1 - \mu}$; differentiating and using $\Gamma (s + 2 - \mu) = (s + 1 - \mu)
  \Gamma (s + 1 - \mu)$ yields \eqref{eq:Dmu_power}.
\end{proof}

\begin{remark}
  [The case $s = 0$]\label{rem:Dmu_const} The hypothesis $s > 0$
  in~\eqref{eq:Dmu_power} cannot be relaxed to $s \ge 0$. For $s = 0$, $f (t)
  \equiv 1$, $f (0) = 1$, and the correction term in~\eqref{eq:caputo} gives
  $D^{\mu} (1) = t^{- \mu} / \Gamma (1 - \mu) - t^{- \mu} / \Gamma (1 - \mu) =
  0$, whereas the right-hand side of~\eqref{eq:Dmu_power} evaluated at $s = 0$
  equals $t^{- \mu} / \Gamma (1 - \mu) \ne 0$. The Caputo derivative of a
  constant is therefore zero, consistent with the convention $D^{\mu} y (0) =
  0$ for solutions of~\eqref{eq:frac_riccati} that vanish at the origin.
\end{remark}

\subsection{The M{\"u}ntz basis}

Fix $\mu \in (0, 1)$ and $T > 0$.

\begin{definition}
  [M{\"u}ntz system]\label{def:muntz} For $k \in \N$, define $e_k (t) \assign
  t^{k \mu}$ on $[0, T]$. The span of $\{1, e_k : k \ge 1\}$ is denoted
  \[ \mathcal{M} \assign \mathrm{span} \{1, t^{\mu}, t^{2 \mu}, \ldots\} . \]
\end{definition}

\begin{remark}
  [Completeness]\label{rem:muntz_complete} By the M{\"u}ntz--Sz{\'a}sz
  theorem, $\{t^{k \mu} \}_{k \ge 1}$ is complete in $L^2 [0, 1]$ iff $\sum_{k
  \ge 1} (k \mu)^{- 1} = \infty$. Since $\sum_{k \ge 1} k^{- 1} = \infty$,
  completeness holds for every $\mu \in (0, 1)$. This underwrites the
  M{\"u}ntz--Tau spectral method.
\end{remark}

\begin{proposition}
  [Closure of $\mathcal{M}$]\label{prop:closure_M} For $k \ge 1$ and $r > 0$,
  
  \begin{align}
    D^{\mu} t^{k \mu} & = \frac{\Gamma (k \mu + 1)}{\Gamma ((k - 1) \mu + 1)} 
    \hspace{0.17em} t^{(k - 1) \mu}  \label{eq:Dmu_Muntz}\\
    I^r t^{k \mu} & = \frac{\Gamma (k \mu + 1)}{\Gamma (k \mu + r + 1)} 
    \hspace{0.17em} t^{k \mu + r}  \label{eq:Ir_Muntz}\\
    t^{j \mu}  \hspace{0.17em} t^{k \mu} & = t^{(j + k) \mu} 
    \label{eq:prod_Muntz}
  \end{align}
  
  Consequently $D^{\mu} \mathcal{M} \subset \mathcal{M}$, $I^r \mathcal{M}
  \subset \mathcal{M}$, and $\mathcal{M}$ is closed under multiplication.
\end{proposition}

\begin{proof}
  \eqref{eq:Dmu_Muntz} follows from  Proposition~\ref{prop:powers} with $s = k
  \mu > 0$; \eqref{eq:Ir_Muntz} from \eqref{eq:Ir_power} with $s = k \mu > -
  1$; \eqref{eq:prod_Muntz} is direct.
\end{proof}

These three closure properties together --- under $D^{\mu}$, under $I^r$, and
under multiplication --- are exactly what makes the M{\"u}ntz lattice the
canonical basis for any constant-coefficient fractional ODE with polynomial
nonlinearity. The M{\"u}ntz--Tau substitution {\cite{esmaeili2015}} preserves
the entire algebraic structure of the equation.

\section{The Constant-Coefficient Fractional Riccati
Equation}\label{sec:riccati}

\subsection{Equation and Volterra representation}

Let $\mu \in (0, 1)$ and $c_0, c_1, c_2 \in \C$. Consider
\begin{equation}
  \label{eq:frac_riccati} D^{\mu} y (t) = c_0 + c_1  \hspace{0.17em} y (t) +
  c_2  \hspace{0.17em} y (t)^2, \qquad I^{1 - \mu} y (0) = 0.
\end{equation}
Applying $I^{\mu}$ and using \eqref{eq:caputo_RL} (and $y (0) = 0$ from 
Remark~\ref{rem:correction}) gives the equivalent Volterra integral equation
\begin{equation}
  \label{eq:frac_riccati_V} y (t) = \frac{1}{\Gamma (\mu)}  \int_0^t (t -
  s)^{\mu - 1}  \hspace{-0.17em} \left( c_0 + c_1  \hspace{0.17em} y (s) + c_2
  \hspace{0.17em} y (s)^2 \right) \dd s.
\end{equation}
\begin{remark}
  [H{\"o}lder regularity at the origin]\label{rem:regularity} Evaluating
  \eqref{eq:frac_riccati_V} at $t = 0$ gives $y (0) = 0$. Near the origin, $y
  (t) = \frac{c_0}{\Gamma (\mu + 1)}  \hspace{0.17em} t^{\mu} +\mathcal{O}
  (t^{2 \mu})$, so $y$ is H{\"o}lder continuous of order $\mu$ but not $C^1$
  at $0$. The H{\"o}lder exponent $\mu$ is the natural regularity scale of the
  problem and the reason ordinary Taylor series in $t$ fail.
\end{remark}

\subsection*{The M{\"u}ntz--Tau Puiseux series solution}

\begin{theorem}
  [Series solution in $\{t^{k \mu} \}$]\label{thm:riccati_series} There exists
  $R > 0$ and coefficients $(a_k)_{k \ge 1} \subset \C$ such that
  \begin{equation}
    \label{eq:y_series} y (t) = \sum_{k = 1}^{\infty} a_k  \hspace{0.17em}
    t^{k \mu}, \qquad |t| < R^{1 / \mu},
  \end{equation}
  with
  
  \begin{align}
    a_1 & = \frac{c_0}{\Gamma (\mu + 1)},  \label{eq:a1_const}\\
    a_{k + 1} & = \frac{\Gamma (k \mu + 1)}{\Gamma ((k + 1) \mu + 1)}  \left(
    c_1  \hspace{0.17em} a_k + c_2  \hspace{-0.17em} \hspace{-0.17em}
    \sum_{\tmscript{\begin{array}{c}
      j + \ell = k\\
      1 \le j, \ell \le k - 1
    \end{array}}} \hspace{-0.17em} \hspace{-0.17em} a_j  \hspace{0.17em}
    a_{\ell} \right), \qquad k \ge 1.  \label{eq:ak_const}
  \end{align}
\end{theorem}

\begin{remark}
  [Quadratic convolution at $k = 1$]\label{rem:conv_k1} For $k = 1$, the
  convolution sum is empty (no pairs $(j, \ell)$ with $j, \ell \ge 1$ and $j +
  \ell = 1$), so $a_2 = \frac{\Gamma (\mu + 1)}{\Gamma (2 \mu + 1)} 
  \hspace{0.17em} c_1  \hspace{0.17em} a_1$. The quadratic term first
  contributes at $k = 2$, giving $a_3 = \frac{\Gamma (2 \mu + 1)}{\Gamma (3
  \mu + 1)}  (c_1  \hspace{0.17em} a_2 + c_2  \hspace{0.17em} a_1^2)$.
\end{remark}

Existence of a positive radius $R > 0$ and analyticity of $y$ in $z = t^{\mu}$
on $|z| < R$ is established in  Theorem~\ref{thm:radius} below by a Banach
fixed-point argument on \eqref{eq:frac_riccati_V}. Granting that, write $y (t)
= \sum_{k \ge 1} a_k  \hspace{0.17em} t^{k \mu}$; by \eqref{eq:prod_Muntz},
\begin{equation}
  y (t)^2 = \sum_{m = 2}^{\infty} \hspace{-0.17em} \left(
  \sum_{\tmscript{\begin{array}{c}
    j + \ell = m\\
    1 \le j, \ell \le m - 1
  \end{array}}} \hspace{-0.17em} a_j  \hspace{0.17em} a_{\ell} \right) t^{m
  \mu}
\end{equation}
abbreviated $b_m \assign \sum_{\tmscript{\begin{array}{c}
  j + \ell = m\\
  1 \le j, \ell \le m - 1
\end{array}}} a_j  \hspace{0.17em} a_{\ell}$ (with $b_1 = 0$). Substituting
into \eqref{eq:frac_riccati_V} and splitting the integral,
\begin{equation}
  \begin{array}{ll}
    y (t) & = \frac{c_0}{\Gamma (\mu)}  \hspace{-0.17em} \int_0^t
    \hspace{-0.17em} (t - s)^{\mu - 1} \dd s + \frac{c_1}{\Gamma (\mu)} 
    \sum_{k \ge 1} a_k  \hspace{-0.17em} \int_0^t \hspace{-0.17em} (t -
    s)^{\mu - 1} s^{k \mu} \dd s\\
    & \quad + \frac{c_2}{\Gamma (\mu)}  \sum_{m \ge 2} b_m  \hspace{-0.17em}
    \int_0^t \hspace{-0.17em} (t - s)^{\mu - 1} s^{m \mu} \dd s
  \end{array}
\end{equation}
The first integral equals $t^{\mu} / \Gamma (\mu + 1)$. For the others,
\eqref{eq:Ir_power} with $s = k \mu$ and $r = \mu$ gives
\[ \int_0^t (t - s)^{\mu - 1} s^{k \mu} \dd s = \Gamma (\mu) \hspace{0.17em}
   \frac{\Gamma (k \mu + 1)}{\Gamma ((k + 1) \mu + 1)}  \hspace{0.17em} t^{(k
   + 1) \mu} \]
Reindexing with $n = k + 1$ and $n = m + 1$ and matching the coefficient of
$t^{n \mu}$ for each $n \ge 1$ yields \eqref{eq:a1_const} for $n = 1$ and
\[ a_n = c_1  \hspace{0.17em} a_{n - 1}  \hspace{0.17em} \frac{\Gamma ((n - 1)
   \mu + 1)}{\Gamma (n \mu + 1)} + c_2  \hspace{0.17em} b_{n - 1} 
   \hspace{0.17em} \frac{\Gamma ((n - 1) \mu + 1)}{\Gamma (n \mu + 1)}, \qquad
   n \ge 2, \]
which is \eqref{eq:ak_const} after renaming $n \mapsto k + 1$.

\begin{theorem}
  [Strictly positive Puiseux radius; coefficient bound; generic
  finiteness]\label{thm:radius} Consider the Volterra
  equation~\eqref{eq:frac_riccati_V}.
  \begin{enumerate}
    \item {\tmem{Existence and $R > 0$.}} Set
    \begin{equation}
      \label{eq:T0_choice} T_0^{\mu} \hspace{0.27em} \assign \hspace{0.27em}
      \min \hspace{-0.17em} \left( \frac{\Gamma (\mu + 1)}{|c_0 | + |c_1 | +
      |c_2 | + 1}, \hspace{0.27em} \frac{\Gamma (\mu + 1)}{2 (|c_1 | + 2| c_2
      | + 1)} \right) \hspace{0.27em} > \hspace{0.27em} 0.
    \end{equation}
    There exists a unique solution $y \in C [0, T_0]$ of
    \eqref{eq:frac_riccati_V} with $\|y\|_{\infty} \le 1$; this solution is
    analytic in $z = t^{\mu}$ on the disc $\{|z| < T_0^{\mu} \}$, so the
    Puiseux series \eqref{eq:y_series} has radius of convergence
    \begin{equation}
      R \hspace{0.27em} \ge \hspace{0.27em} T_0^{\mu} \hspace{0.27em} >
      \hspace{0.27em} 0
    \end{equation}
    \item {\tmem{Coefficient bound.}} For every $0 < r < R$, with $M (r)
    \assign \sup_{|z| = r} |g (z) |$ where $g (z) \assign y (z^{1 / \mu})$,
    \begin{equation}
      \label{eq:cauchy_bound} |a_k | \hspace{0.27em} \le \hspace{0.27em}
      \frac{M (r)}{r^k} \qquad \forall k \ge 1
    \end{equation}
    \item {\tmem{Generic finiteness, $R < \infty$.}} For generic parameters
    with $c_2 \ne 0$, $R < \infty$: the complexified Riccati flow develops
    movable singularities at finite distance from the origin in the complex
    $z$-plane.
  \end{enumerate}
\end{theorem}

\begin{proof}
  \begin{enumeratenumeric}
    \item Let $B_1 \assign \{y \in C [0, T_0] : \|y\|_{\infty} \le 1\}$ with
    the supremum norm, and define
    \begin{equation}
      \Phi (y) (t) \hspace{0.27em} \assign \hspace{0.27em} \frac{1}{\Gamma
      (\mu)}  \int_0^t (t - s)^{\mu - 1}  \left( c_0 + c_1  \hspace{0.17em} y
      (s) + c_2  \hspace{0.17em} y (s)^2 \right) \dd s \hspace{0.27em} =
      \hspace{0.27em} I^{\mu}  (c_0 + c_1 y + c_2 y^2) (t)
    \end{equation}
    For $y \in B_1$ and $t \in [0, T_0]$, $|c_0 + c_1 y (t) + c_2 y (t)^2 |
    \le |c_0 | + |c_1 | + |c_2 |$, so by \eqref{eq:Ir_power},
    \begin{equation}
      | \Phi (y) (t) | \hspace{0.27em} \le \hspace{0.27em} (|c_0 | + |c_1 | +
      |c_2 |)  \frac{t^{\mu}}{\Gamma (\mu + 1)} \hspace{0.27em} \le
      \hspace{0.27em} (|c_0 | + |c_1 | + |c_2 |)  \frac{T_0^{\mu}}{\Gamma (\mu
      + 1)}
    \end{equation}
    By the first term in \eqref{eq:T0_choice},
    \begin{equation}
      \frac{T_0^{\mu}  (|c_0 | + |c_1 | + |c_2 |)}{\Gamma (\mu + 1)} \le
      \frac{|c_0 | + |c_1 | + |c_2 |}{|c_0 | + |c_1 | + |c_2 | + 1} < 1
    \end{equation}
    ; hence $\Phi (B_1) \subset B_1$. For $y_1, y_2 \in B_1$,
    \begin{equation}
      | c_1 (y_1 - y_2) + c_2 (y_1^2 - y_2^2) | \hspace{0.27em} =
      \hspace{0.27em} | c_1 + c_2 (y_1 + y_2) |  \hspace{0.17em} |y_1 - y_2 |
      \hspace{0.27em} \le \hspace{0.27em} (|c_1 | + 2| c_2 |)  \hspace{0.17em}
      \|y_1 - y_2 \|_{\infty}
    \end{equation}
    so
    \begin{equation}
      \| \Phi (y_1) - \Phi (y_2)\|_{\infty} \hspace{0.27em} \le
      \hspace{0.27em} \frac{T_0^{\mu}}{\Gamma (\mu + 1)} (|c_1 | + 2| c_2 |) 
      \hspace{0.17em} \|y_1 - y_2 \|_{\infty}
    \end{equation}
    By the second term in \eqref{eq:T0_choice}, $T_0^{\mu} (|c_1 | + 2| c_2 |)
    / \Gamma (\mu + 1) \le (|c_1 | + 2| c_2 |) / [2 (|c_1 | + 2| c_2 | + 1)] <
    \tfrac{1}{2}$; hence $\Phi$ is a strict contraction on $B_1$. The Banach
    fixed-point theorem provides a unique $y^{\ast} \in B_1$ with $\Phi
    (y^{\ast}) = y^{\ast}$, which solves \eqref{eq:frac_riccati_V}. The formal
    Puiseux series defined by \eqref{eq:a1_const}--\eqref{eq:ak_const} solves
    the same Volterra equation order-by-order; by uniqueness, the series
    equals $y^{\ast}$ on $[0, T_0]$, and hence converges for $|t^{\mu} | <
    T_0^{\mu}$, proving $R \ge T_0^{\mu} > 0$.
    
    \item By item~(i), $g (z) = y (z^{1 / \mu}) = \sum_{k \ge 1} a_k z^k$ is
    analytic on $\{|z| < R\}$. For any $0 < r < R$, Cauchy's integral formula
    gives
    \begin{equation}
      a_k \hspace{0.27em} = \hspace{0.27em} \frac{1}{2 \pi \ii}  \oint_{|z| =
      r} \frac{g (z)}{z^{k + 1}} \hspace{0.17em} \dd z
    \end{equation}
    \begin{equation}
      |a_k | \hspace{0.27em} \le \hspace{0.27em} \sup_{|z| = r} |g (z) | \cdot
      r^{- k} \hspace{0.27em} = \hspace{0.27em} M (r)  \hspace{0.17em} r^{- k}
    \end{equation}
    \item For the classical case $\mu = 1$ with $c_2 \ne 0$, the linear
    substitution $y = - w' / (c_2 w)$ reduces $y' = c_0 + c_1 y + c_2 y^2$ to
    a second-order linear ODE with constant coefficients for $w$; the zeros of
    $w$ generically lie at finite $t$, and at each such zero $y$ has a simple
    pole. The same algebraic structure of the recurrence~\eqref{eq:ak_const}
    persists for $\mu \in (0, 1)$, and the resulting $g (z)$ generically
    inherits movable singularities at finite distance from the origin in the
    $z$-plane (the exact location of the nearest singularity is encoded by the
    limit of the diagonal Pad{\'e} poles, see  Section~\ref{sec:pade}). Hence
    $R < \infty$ for generic parameters with $c_2 \ne 0$
  \end{enumeratenumeric}
\end{proof}

\begin{remark}
  [Why this is only a local representation]\label{rem:radius_local}
  Theorem~\ref{thm:radius}(iii) is the central obstruction: even though every
  $a_k$ is computed in closed form by \eqref{eq:ak_const}, the series
  \eqref{eq:y_series} does {\tmem{not}} converge for all $t \ge 0$. The
  complex $z$-plane carries movable singularities of the Riccati flow, and the
  Puiseux radius $R$ is the distance to the nearest such singularity. On the
  positive real $t$-axis the solution $y (t)$ may exist for all $t$, yet the
  M{\"u}ntz series ceases to converge once $t^{\mu}$ exceeds $R$. A genuine
  global representation must exploit {\tmem{analytic continuation}}, not raw
  series summation. This is the role played by Pad{\'e} approximation in the
  next section.
\end{remark}

\section{Pad{\'e} Resummation in $z = t^{\mu}$}\label{sec:pade}

\subsection{Power series in $z$ and the Hankel system}

Let $z \assign t^{\mu}$ and define $g (z) \assign y (z^{1 / \mu})$. Then
\begin{equation}
  \label{eq:g_series} g (z) = \sum_{k = 1}^{\infty} a_k  \hspace{0.17em} z^k,
  \qquad |z| < R.
\end{equation}
Fix $M \in \N$ and seek polynomials
\begin{equation}
  P_M (z) \assign \sum_{k = 0}^M p_k  \hspace{0.17em} z^k, \qquad Q_M (z)
  \assign 1 + \sum_{k = 1}^M q_k  \hspace{0.17em} z^k
\end{equation}
satisfying the Pad{\'e} matching condition
\begin{equation}
  \label{eq:Pade_match} Q_M (z)  \hspace{0.17em} g (z) - P_M (z) =\mathcal{O}
  (z^{2 M + 1})  \quad (z \to 0) .
\end{equation}
\begin{lemma}
  [Hankel form of the denominator system]\label{lem:hankel} Substituting
  \eqref{eq:g_series} into \eqref{eq:Pade_match} and matching coefficients of
  $z^n$ for $0 \le n \le 2 M$ gives $p_0 = 0$,
  
  \begin{align}
    p_n & = a_n + \sum_{j = 1}^{\min (n, M)} q_j  \hspace{0.17em} a_{n - j},
    \qquad 1 \le n \le M  \label{eq:Pade1}\\
    0 & = a_n + \sum_{j = 1}^M q_j  \hspace{0.17em} a_{n - j}, \qquad M + 1
    \le n \le 2 M  \label{eq:Pade2}
  \end{align}
  
  The denominator system \eqref{eq:Pade2} is the Hankel linear system
  \begin{equation}
    \label{eq:hankel_system} H_M  \hspace{0.17em} \tmmathbf{q}= -\tmmathbf{b},
    \qquad H_M = \left(\begin{array}{cccc}
      a_M & a_{M - 1} & \cdots & a_1\\
      a_{M + 1} & a_M & \cdots & a_2\\
      \vdots & \vdots & \ddots & \vdots\\
      a_{2 M - 1} & a_{2 M - 2} & \cdots & a_M
    \end{array}\right), \quad \tmmathbf{q}= \left(\begin{array}{c}
      q_1\\
      \vdots\\
      q_M
    \end{array}\right), \quad \tmmathbf{b}= \left(\begin{array}{c}
      a_{M + 1}\\
      \vdots\\
      a_{2 M}
    \end{array}\right) .
  \end{equation}
\end{lemma}

\begin{proof}
  Direct expansion of $Q_M  \hspace{0.17em} g$ and identification of
  coefficients.
\end{proof}

\begin{remark}
  [Invertibility of $H_M$]\label{rem:HM_invertible} $H_M$ is invertible iff $g
  (z)$ is not a rational function of degree less than $M$. For $\mu = 1$
  (classical case) the Riccati solution is rational, and $H_M$ becomes
  singular for $M$ beyond the exact pole order. For $\mu \in (0, 1)$, $g (z)$
  admits no finite-order rational representation; $\det H_M \ne 0$ generically
  for all $M \ge 1$. If a particular order $M$ yields $\det H_M = 0$, use $M
  \pm 1$.
\end{remark}

\begin{lemma}
  [Solution]\label{lem:Pade_system} If $\det H_M \neq 0$, then
  \eqref{eq:hankel_system} has the unique solution $\tmmathbf{q}= - H_M^{- 1}
  \tmmathbf{b}$, and the numerator coefficients are given by \eqref{eq:Pade1}.
\end{lemma}

The diagonal Pad{\'e} approximant is then
\begin{equation}
  \label{eq:R_M_def} R_M (z) \hspace{0.27em} \assign \hspace{0.27em} \frac{P_M
  (z)}{Q_M (z)} \hspace{0.27em} = \hspace{0.27em} \frac{\sum_{k = 0}^M p_k 
  \hspace{0.17em} z^k}{1 + \sum_{k = 1}^M q_k  \hspace{0.17em} z^k} .
\end{equation}
\begin{definition}
  [Pad{\'e} domain]\label{def:D_M} $\mathcal{D}_M \assign \{z \in \C : \det
  H_M \neq 0, Q_M (z) \neq 0\}$.
\end{definition}

\subsection{Remainder representation and successive-difference bound}

Define the remainder
\begin{equation}
  \label{eq:eps_g} \varepsilon_M (z) \assign g (z) - R_M (z)
\end{equation}
By construction $\varepsilon_M (z) =\mathcal{O} (z^{2 M + 1})$, and standard
Pad{\'e} theory {\cite{baker1996}} gives the structural representation
\begin{equation}
  \label{eq:eps_structure} \varepsilon_M (z) = \frac{\Pi_{M + 1} (z)}{Q_M
  (z)^2}  \hspace{0.17em} z^{2 M + 1}
\end{equation}
for some polynomial $\Pi_{M + 1}$ of degree at most $M + 1$.

For the a-posteriori bound, define
\begin{equation}
  \label{eq:Delta_def} \Delta_k (z) \assign R_k (z) - R_{k - 1} (z), \qquad k
  \ge 1.
\end{equation}
\begin{assumption}
  [Uniform geometric tail]\label{ass:diff_bound} At a fixed $z \in
  \mathcal{D}_M \cap \mathcal{D}_{M - 1}$:
  \begin{enumerate}
    \item $g$ has no poles in $\{| \zeta | \le |z|\}$;
    
    \item $Q_M (z) \ne 0$ and $Q_{M - 1} (z) \ne 0$;
    
    \item there exists $\rho \in (0, 1)$ such that $| \Delta_{k + 1} (z) | \le
    \rho \hspace{0.17em} | \Delta_k (z) |$ for every $k \ge M - 1$.
  \end{enumerate}
\end{assumption}

\begin{theorem}
  [Computable error bound]\label{thm:error} Under ~\ref{ass:diff_bound} with
  contraction ratio $\rho$, the Pad{\'e} sequence $R_k (z)$ converges and
  \begin{equation}
    \label{eq:error_bound_uniform} |g (z) - R_M (z) | \hspace{0.27em} \le
    \hspace{0.27em} \frac{\rho \hspace{0.17em} | \Delta_M (z) |}{1 - \rho} .
  \end{equation}
  In particular, taking the smallest a-posteriori-observable contraction ratio
  $\rho_{\ast} \assign | \Delta_M (z) | / | \Delta_{M - 1} (z) |$ (which
  equals $\rho$ at the index $k = M - 1$ by~(iii) with equality, hence
  satisfies $\rho_{\ast} \le \rho < 1$), and assuming~(iii) holds with this
  $\rho_{\ast}$,
  \begin{equation}
    \label{eq:error_bound} |g (z) - R_M (z) | \hspace{0.27em} \le
    \hspace{0.27em} \frac{| \Delta_M (z) |^2}{| \Delta_{M - 1} (z) | - |
    \Delta_M (z) |} .
  \end{equation}
\end{theorem}

\begin{proof}
  By ~\ref{ass:diff_bound}(iii), $| \Delta_{M + j} (z) | \le \rho^j | \Delta_M
  (z) |$ for every $j \ge 0$ (induction on $j$, with the base case from~(iii)
  and the step from $| \Delta_{M + j + 1} | \le \rho | \Delta_{M + j} |$).
  Since $\sum_{j = 0}^{\infty} \rho^j | \Delta_M (z) | < \infty$, the
  telescoping series
  \begin{equation}
    g (z) - R_M (z) \hspace{0.27em} = \hspace{0.27em} \sum_{k = M +
    1}^{\infty} \Delta_k (z)
  \end{equation}
  converges absolutely. Estimating term-by-term,
  \begin{equation}
    |g (z) - R_M (z) | \hspace{0.27em} \le \hspace{0.27em} \sum_{j =
    1}^{\infty} | \Delta_{M + j} (z) | \hspace{0.27em} \le \hspace{0.27em} |
    \Delta_M (z) |  \sum_{j = 1}^{\infty} \rho^j \hspace{0.27em} =
    \hspace{0.27em} \frac{\rho \hspace{0.17em} | \Delta_M (z) |}{1 - \rho}
  \end{equation}
  which is~\eqref{eq:error_bound_uniform}. Substituting $\rho = \rho_{\ast} =
  | \Delta_M (z) | / | \Delta_{M - 1} (z) |$ yields $\rho_{\ast} | \Delta_M |
  / (1 - \rho_{\ast}) = | \Delta_M |^2 / (| \Delta_{M - 1} | - | \Delta_M |)$,
  giving~\eqref{eq:error_bound}.
\end{proof}

\subsection{Global convergence on the positive real axis}

\begin{theorem}
  [de Montessus de Ballore {\cite{baker1996,montessus1902}}]\label{thm:dMdB}
  Let $g$ be meromorphic in $\{|z| < \rho\}$ with exactly $n$ poles (counted
  with multiplicity) in that disc and no poles on the boundary. The diagonal
  Pad{\'e} approximants $R_M$ of order $(M, M)$ converge to $g$ uniformly on
  compact subsets of $\{|z| < \rho\} \setminus \{ \text{poles} \}$ as $M \to
  \infty$.
\end{theorem}

\begin{theorem}
  [Nuttall--Pommerenke {\cite{baker1996,nuttall1970}}]\label{thm:NP} Let $g$
  be meromorphic in $\C$ with at most countably many poles having no finite
  accumulation point. The diagonal sequence $R_M$ converges to $g$ {\tmem{in
  capacity}} on $\C \setminus E$, where $E$ has logarithmic capacity zero. In
  particular, convergence is quasi-everywhere, though not necessarily
  pointwise at every individual point.
\end{theorem}

\begin{remark}
  [Which theorem applies]\label{rem:which_thm} Theorem~\ref{thm:dMdB} gives
  {\tmem{pointwise}} uniform convergence on compacta but requires knowing the
  exact pole count.  Theorem~\ref{thm:NP} requires no such count but yields
  only quasi-everywhere convergence, and presupposes meromorphy on all of
  $\C$. When $g (z) = y (z^{1 / \mu})$ is meromorphic in $\C$ with no poles on
  $(0, \infty)$ --- the canonical case for globally bounded Riccati solutions
  on the real time axis --- both theorems apply, and $R_M (t^{\mu}) \to y (t)$
  for every $t > 0$.
\end{remark}

\begin{theorem}
  [Global Pad{\'e}--M{\"u}ntz representation]\label{thm:pade_riccati} Let $y
  (t)$ solve \eqref{eq:frac_riccati}, assume $y$ is bounded on $[0, \infty)$
  (so that $g (z)$ has no poles on the positive real $z$-axis), and assume
  that $g (z) \assign y (z^{1 / \mu})$ admits a meromorphic continuation to
  $\C$ with at most countably many poles having no finite accumulation point.
  Then
  \begin{equation}
    \label{eq:y_global_pade} \hspace{0.17em} y (t) = \lim_{M \to \infty} R_M
    (t^{\mu}) = \lim_{M \to \infty}  \frac{P_M (t^{\mu})}{Q_M (t^{\mu})} 
    \quad \text{{\forall}} t \ge 0
  \end{equation}
\end{theorem}

\begin{proof}
  Under the meromorphy hypothesis,  Theorem~\ref{thm:NP} applies and $R_M \to
  g$ in capacity on $\C \setminus E$ with $\mathrm{cap} (E) = 0$. Since $g$ is
  holomorphic on a neighbourhood of every point of $(0, \infty)$ (by the
  no-positive-pole hypothesis), and meromorphic on every disc $\{|z| < \rho\}$
  with finitely many poles in that disc,  Theorem~\ref{thm:dMdB} gives uniform
  convergence on compact subsets of $\C \setminus \{ \text{poles} \}$; in
  particular pointwise convergence at every $z = t^{\mu} \in (0, \infty)$. At
  $t = 0$, $y (0) = 0$ and $R_M (0) = p_0 / Q_M (0) = 0 / 1 = 0$ for every
  $M$, so the limit holds trivially.
\end{proof}

\begin{remark}
  [Status of the meromorphy hypothesis]\label{rem:meromorphy} Meromorphy of
  $g$ on $\C$ is established for $\mu = 1$ (where the linearization $y = - w'
  / (c_2 w)$ reduces the Riccati equation to a linear ODE with entire
  solutions) and is conjectured for $\mu \in (0, 1)$ on the basis of the
  Painlev{\'e} property of the underlying flow and the empirical pole
  distribution of the diagonal Pad{\'e} sequence. A complete proof for
  arbitrary $\mu \in (0, 1)$ lies outside the scope of the present paper and
  is connected to the Riemann--Hilbert structure discussed in 
  Section~\ref{sec:resurgence}. In practical computation, the hypothesis is
  verified a~posteriori by the convergence behavior of $R_M (t^{\mu})$ at
  probe points $t$ of interest, controlled by the
  bound~\eqref{eq:error_bound}.
\end{remark}

\subsection{Algorithmic summary}

Given $(c_0, c_1, c_2, \mu)$ and Pad{\'e} order $M$:
\begin{enumerate}
  \item Compute the Puiseux coefficients $a_1, \ldots, a_{2 M}$ by the
  recurrence \eqref{eq:a1_const}--\eqref{eq:ak_const}. Cost: $\mathcal{O}
  (M^2)$.
  
  \item Form the Hankel matrix $H_M$ and right-hand side $\tmmathbf{b}$ from
  \eqref{eq:hankel_system}.
  
  \item Solve $H_M \tmmathbf{q}= -\tmmathbf{b}$ for the denominator
  coefficients. Cost: $\mathcal{O} (M^3)$.
  
  \item Recover the numerator coefficients via \eqref{eq:Pade1}.
  
  \item Evaluate $y_M (t) = P_M (t^{\mu}) / Q_M (t^{\mu})$ at any $t \ge 0$.
  Cost: $\mathcal{O} (M)$ per query.
\end{enumerate}
Steps 1--4 are independent of $t$; the entire family $\{y_M (t) : t \ge 0\}$
is encoded in a single rational function in $z = t^{\mu}$. There is no grid,
no time stepping, no Newton iteration, no fixed-point loop.

\

\begin{thebibliography}{9}
  {\bibitem{eleuch2019}}O.~El~Euch and M.~Rosenbaum, {\tmem{The characteristic
  function of rough Heston models}}, Mathematical Finance, 29(1):3--38, 2019.
  
  {\bibitem{baker1996}}G.{\hspace{0.17em}}A.~Baker and P.~Graves-Morris,
  {\tmem{Pad{\'e} Approximants}}, 2nd ed., Cambridge University Press, 1996.
  
  {\bibitem{esmaeili2015}}S.~Esmaeili and M.~Shamsi, {\tmem{The
  M{\"u}ntz--Legendre Tau method for fractional differential equations}},
  Applied Mathematical Modelling, 40(2):671--684, 2015.
  
  {\bibitem{montessus1902}}R.~de~Montessus de~Ballore, {\tmem{Sur les
  fractions continues alg{\'e}briques}}, Bull. Soc. Math. France, 30:28--36,
  1902.
  
  {\bibitem{nuttall1970}}J.~Nuttall, {\tmem{The convergence of Pad{\'e}
  approximants of meromorphic functions}}, J. Math. Anal. Appl.,
  31(1):147--153, 1970.
  
  {\bibitem{pommerenke1973}}Ch.~Pommerenke, {\tmem{Pad{\'e} approximants and
  convergence in capacity}}, J. Math. Anal. Appl., 41:775--780, 1973.
  
  {\bibitem{stahl1997}}H.~Stahl, {\tmem{The convergence of Pad{\'e}
  approximants to functions with branch points}}, J. Approx. Theory,
  91:139--204, 1997.
  
  {\bibitem{ecalle1981}}J.~{\'E}calle, {\tmem{Les fonctions r{\'e}surgentes}},
  Vols. I--III, Publ. Math. d'Orsay, 1981--1985.
  
  {\bibitem{costin2008}}O.~Costin, {\tmem{Asymptotics and Borel Summability}},
  Chapman \& Hall/CRC, 2008.
  
  {\bibitem{deift1999}}P.~Deift, {\tmem{Orthogonal Polynomials and Random
  Matrices: A Riemann--Hilbert Approach}}, Courant Lecture Notes 3, AMS, 1999.
  
  {\bibitem{callegaro2021}}G.~Callegaro, M.~Grasselli, and G.~Pag{\`e}s,
  {\tmem{Fast hybrid schemes for fractional Riccati equations (rough is not so
  tough)}}, Mathematics of Operations Research, 46(1):221--254, 2021.
  
  {\bibitem{gatheral2019}}J.~Gatheral and R.~Radoi{\v c}i{\'c},
  {\tmem{Rational approximation of the rough Heston solution}}, International
  Journal of Theoretical and Applied Finance, 22(3):1950010, 2019.
\end{thebibliography}

\end{document}
